% =================================================
% 文件: LDAP.tex
% 章节: LDAP 轻量级目录访问协议入门
% 版本: v1.0
% =================================================

\documentclass[12pt,UTF8]{ctexbook}

\input{../../../common/page}

\xeCJKsetup{AutoFallBack=true}
\setCJKfamilyfont{hei}{SimHei}
\setCJKfamilyfont{kai}{KaiTi}

\usepackage{titletoc}
\titlecontents{chapter}[0pt]{\vspace{3mm}\bf\addvspace{2pt}\filright}{\contentspush{\thecontentslabel\hspace{0.8em}}}{}{\titlerule*[8pt]{.}\contentspage}

\usepackage[colorlinks,linkcolor=blue,citecolor=blue,urlcolor=blue]{hyperref}

\ctexset{
	part/name= {第,卷},
	part/number={\chinese{part}},
	chapter/name={第,章},
	chapter/number={\arabic{chapter}}
}

\usepackage{graphicx}
\graphicspath{{Images/}}

\usepackage[perpage]{footmisc}

\usepackage{enumitem}

\title{\heiti\zihao{0} LDAP 轻量级目录访问协议入门指南}
\author{WangFei}
\date{\today}

\begin{document}

\maketitle
\tableofcontents

\frontmatter
\chapter{前言}

本指南面向初学者，详细介绍 LDAP（Lightweight Directory Access Protocol）轻量级目录访问协议的原理、配置和使用方法。

\mainmatter

\chapter{LDAP 基础概念}
\label{chap:ldap_basics}

\section{什么是 LDAP}
\label{sec:what_is_ldap}

LDAP（Lightweight Directory Access Protocol，轻量级目录访问协议）是一种用于访问和维护目录服务的协议。目录服务是一种特殊的数据库，用于存储和管理层次化的信息，如用户信息、组织架构、设备信息等。

LDAP 的主要特点：
\begin{itemize}
    \item \textbf{轻量级}：基于 TCP/IP，协议简单高效
    \item \textbf{层次化结构}：数据以树形结构组织（LDAP 目录树）
    \item \textbf{只读优化}：适合查询操作，写入性能相对较低
    \item \textbf{分布式}：支持多服务器复制和分布式查询
    \item \textbf{标准协议}：被广泛应用于企业身份认证和目录服务
\end{itemize}

为什么需要 LDAP？
\begin{itemize}
    \item \textbf{集中管理}：统一管理用户、组、设备等信息
    \item \textbf{单点登录}：实现企业级单点登录（SSO）
    \item \textbf{身份认证}：作为统一的身份认证源
    \item \textbf{数据共享}：在多个应用之间共享目录信息
\end{itemize}

\section{LDAP 目录树结构}
\label{sec:ldap_directory_tree}

LDAP 数据以树形结构组织，称为目录树（Directory Tree）。

基本概念：
\begin{itemize}
    \item \textbf{条目（Entry）}：目录树中的一个节点，包含一组属性
    \item \textbf{属性（Attribute）}：条目的特性，如姓名、电话、邮箱等
    \item \textbf{对象类（Object Class）}：定义条目必须包含的属性
    \item \textbf{可分辨名称（DN）}：条目的唯一标识，如 "cn=张三,ou=users,dc=example,dc=com"
    \item \textbf{相对可分辨名称（RDN）}：DN 的一部分，如 "cn=张三"
\end{itemize}

目录树示例：
\begin{verbatim}
dc=example,dc=com
├── ou=users
│   ├── cn=张三
│   ├── cn=李四
│   └── cn=王五
├── ou=groups
│   ├── cn=开发组
│   └── cn=运维组
└── ou=devices
    └── cn=server01
\end{verbatim}

\section{LDAP 协议版本}
\label{sec:ldap_versions}

LDAP 有三个主要版本：

\begin{table}[htbp]
    \centering
    \caption{LDAP 协议版本}
    \begin{tabular}{|c|p{8cm}|}
        \hline
        \textbf{版本} & \textbf{说明} \\
        \hline
        LDAP v2 & 最早的版本，功能有限，已被弃用 \\
        \hline
        LDAP v3 & 当前主流版本，支持 TLS/SSL、SASL 认证、扩展操作等 \\
        \hline
        LDAP v3 (draft) & 正在开发中的新版本，支持更多高级功能 \\
        \hline
    \end{tabular}
\end{table}

当前所有主流 LDAP 服务器都支持 LDAP v3。

\chapter{LDAP 服务器安装}
\label{chap:ldap_install}

\section{安装 OpenLDAP}
\label{sec:install_openldap}

OpenLDAP 是最流行的开源 LDAP 服务器实现。

\subsection{在 CentOS/RHEL 上安装 OpenLDAP}

\begin{verbatim}
# 使用 dnf 安装
sudo dnf install -y openldap openldap-servers openldap-clients

# 启动服务
sudo systemctl start slapd
sudo systemctl enable slapd

# 允许 LDAP 服务通过防火墙
sudo firewall-cmd --add-service=ldap --permanent
sudo firewall-cmd --add-service=ldaps --permanent
sudo firewall-cmd --reload
\end{verbatim}

\subsection{在 Ubuntu/Debian 上安装 OpenLDAP}

\begin{verbatim}
# 使用 apt 安装
sudo apt update
sudo apt install -y slapd ldap-utils

# 安装过程中会提示设置管理员密码
\end{verbatim}

\section{安装 389 Directory Server}
\label{sec:install_389ds}

389 Directory Server 是 Red Hat 开发的企业级 LDAP 服务器。

\subsection{在 CentOS/RHEL 上安装 389 Directory Server}

\begin{verbatim}
# 安装 389 Directory Server
sudo dnf install -y 389-ds-base

# 创建实例
sudo dscreate create-template /tmp/ds-template.inf

# 修改模板配置
sudo vi /tmp/ds-template.inf

# 根据模板创建实例
sudo dscreate from-file /tmp/ds-template.inf

# 启动服务
sudo systemctl start dirsrv@instance_name
sudo systemctl enable dirsrv@instance_name
\end{verbatim}

\section{验证安装}
\label{sec:verify_install}

安装完成后，可以通过以下命令验证 LDAP 是否正常运行：

\begin{verbatim}
# 检查 slapd 服务状态（OpenLDAP）
sudo systemctl status slapd

# 测试 LDAP 连接
ldapsearch -x -b '' -s base '(objectClass=*)' namingContexts

# 查看服务器信息
ldapsearch -x -b 'cn=config' '(objectClass=*)'
\end{verbatim}

\chapter{LDAP 配置文件详解}
\label{chap:ldap_config}

\section{配置文件结构}
\label{sec:config_structure}

OpenLDAP 的配置文件通常位于以下位置：

\begin{itemize}
    \item \textbf{slapd.conf}：传统配置文件（已弃用）
    \item \textbf{cn=config}：新的配置方式，通过 LDAP 协议管理
    \item \textbf{/etc/openldap/}：配置文件目录
    \item \textbf{/var/lib/openldap/}：数据存储目录
\end{itemize}

\section{使用 cn=config 配置}
\label{sec:cn_config}

OpenLDAP 推荐使用 cn=config 进行配置，这是一种通过 LDAP 协议管理配置的方式。

查看当前配置：
\begin{verbatim}
ldapsearch -x -b 'cn=config' '(objectClass=*)'
\end{verbatim}

修改配置示例：
\begin{verbatim}
# 创建修改文件
cat > /tmp/config.ldif <<EOF
dn: cn=config
changetype: modify
replace: olcLogLevel
olcLogLevel: stats
EOF

# 应用配置修改
ldapmodify -x -D 'cn=admin,cn=config' -W -f /tmp/config.ldif
\end{verbatim}

\section{常用配置选项}
\label{sec:common_options}

常用的 OpenLDAP 配置选项：

\begin{itemize}
    \item \textbf{olcLogLevel}：日志级别
    \item \textbf{olcTLSCertificateFile}：TLS 证书文件
    \item \textbf{olcTLSCertificateKeyFile}：TLS 密钥文件
    \item \textbf{olcTLSCACertificateFile}：CA 证书文件
    \item \textbf{olcAccess}：访问控制规则
    \item \textbf{olcSizeLimit}：查询结果大小限制
    \item \textbf{olcTimeLimit}：查询时间限制
\end{itemize}

\chapter{LDAP 目录结构设计}
\label{chap:ldap_structure}

\section{设计原则}
\label{sec:design_principles}

设计 LDAP 目录结构时应遵循以下原则：

\begin{itemize}
    \item \textbf{层次清晰}：按照组织架构设计目录树
    \item \textbf{易于扩展}：考虑未来的业务增长
    \item \textbf{符合标准}：使用标准的对象类和属性
    \item \textbf{便于管理}：便于权限管理和维护
\end{itemize}

\section{常见目录结构}
\label{sec:common_structures}

常见的 LDAP 目录结构：

\begin{verbatim}
dc=example,dc=com
├── ou=users              # 用户组织单元
│   ├── cn=张三
│   ├── cn=李四
│   └── cn=王五
├── ou=groups             # 组组织单元
│   ├── cn=开发组
│   └── cn=运维组
├── ou=devices            # 设备组织单元
│   └── cn=server01
├── ou=applications       # 应用组织单元
│   └── cn=gitlab
└── ou=services           # 服务组织单元
    └── cn=ldap
\end{verbatim}

\section{对象类和属性}
\label{sec:object_classes}

常用的对象类：

\begin{table}[htbp]
    \centering
    \caption{常用 LDAP 对象类}
    \begin{tabular}{|c|p{8cm}|}
        \hline
        \textbf{对象类} & \textbf{说明} \\
        \hline
        inetOrgPerson & 互联网组织人员，包含姓名、电话、邮箱等属性 \\
        \hline
        organizationalPerson & 组织人员，包含职位等属性 \\
        \hline
        person & 基本人员对象类 \\
        \hline
        organizationalUnit & 组织单元 \\
        \hline
        groupOfNames & 名称组，包含成员列表 \\
        \hline
        groupOfUniqueNames & 唯一名称组 \\
        \hline
        device & 设备对象类 \\
        \hline
    \end{tabular}
\end{table}

常用属性：
\begin{itemize}
    \item \textbf{cn}：通用名称（Common Name）
    \item \textbf{sn}：姓氏（Surname）
    \item \textbf{givenName}：名字
    \item \textbf{mail}：邮箱地址
    \item \textbf{telephoneNumber}：电话号码
    \item \textbf{uid}：用户ID
    \item \textbf{userPassword}：用户密码
    \item \textbf{member}：组成员
\end{itemize}

\chapter{LDAP 基本操作}
\label{chap:ldap_operations}

\section{添加条目}
\label{sec:add_entry}

使用 ldapadd 命令添加条目：

\begin{verbatim}
# 创建 LDIF 文件
cat > /tmp/add_user.ldif <<EOF
dn: cn=张三,ou=users,dc=example,dc=com
objectClass: inetOrgPerson
cn: 张三
sn: 张
givenName: 三
mail: zhangsan@example.com
telephoneNumber: 13800138000
userPassword: password123
EOF

# 添加条目
ldapadd -x -D 'cn=admin,dc=example,dc=com' -W -f /tmp/add_user.ldif
\end{verbatim}

\section{查询条目}
\label{sec:search_entry}

使用 ldapsearch 命令查询条目：

\begin{verbatim}
# 查询所有用户
ldapsearch -x -b 'ou=users,dc=example,dc=com' '(objectClass=inetOrgPerson)'

# 查询特定用户
ldapsearch -x -b 'ou=users,dc=example,dc=com' '(cn=张三)'

# 查询用户的邮箱属性
ldapsearch -x -b 'ou=users,dc=example,dc=com' '(cn=张三)' mail

# 递归查询所有条目
ldapsearch -x -b 'dc=example,dc=com' '(objectClass=*)'
\end{verbatim}

\section{修改条目}
\label{sec:modify_entry}

使用 ldapmodify 命令修改条目：

\begin{verbatim}
# 创建修改文件
cat > /tmp/modify_user.ldif <<EOF
dn: cn=张三,ou=users,dc=example,dc=com
changetype: modify
replace: telephoneNumber
telephoneNumber: 13900139000
EOF

# 修改条目
ldapmodify -x -D 'cn=admin,dc=example,dc=com' -W -f /tmp/modify_user.ldif
\end{verbatim}

修改操作类型：
\begin{itemize}
    \item \textbf{replace}：替换属性值
    \item \textbf{add}：添加属性值
    \item \textbf{delete}：删除属性值
\end{itemize}

\section{删除条目}
\label{sec:delete_entry}

使用 ldapdelete 命令删除条目：

\begin{verbatim}
# 删除单个条目
ldapdelete -x -D 'cn=admin,dc=example,dc=com' -W 'cn=张三,ou=users,dc=example,dc=com'

# 删除整个组织单元（需要先删除子条目）
ldapdelete -x -D 'cn=admin,dc=example,dc=com' -W 'ou=users,dc=example,dc=com'
\end{verbatim}

\chapter{LDAP 访问控制}
\label{chap:ldap_access_control}

\section{访问控制列表（ACL）}
\label{sec:acl}

LDAP 使用访问控制列表（ACL）来控制对目录条目的访问权限。

ACL 配置示例：
\begin{verbatim}
# 创建 ACL 配置文件
cat > /tmp/acl.ldif <<EOF
dn: olcDatabase={1}mdb,cn=config
changetype: modify
replace: olcAccess
olcAccess: to attrs=userPassword by self write by anonymous auth by * none
olcAccess: to dn.base="" by * read
olcAccess: to * by self write by users read by anonymous auth
EOF

# 应用 ACL 配置
ldapmodify -x -D 'cn=admin,cn=config' -W -f /tmp/acl.ldif
\end{verbatim}

ACL 规则说明：
\begin{itemize}
    \item \textbf{to attrs=userPassword}：针对 userPassword 属性
    \item \textbf{by self write}：允许自己写入
    \item \textbf{by anonymous auth}：允许匿名认证
    \item \textbf{by * none}：其他人无权限
    \item \textbf{to dn.base=""}：针对根 DN
    \item \textbf{by * read}：允许所有人读取
\end{itemize}

\section{常见权限配置}
\label{sec:common_permissions}

常见的权限配置场景：

\begin{verbatim}
# 允许用户修改自己的密码
olcAccess: to attrs=userPassword by self write by * none

# 允许用户读取自己的信息
olcAccess: to dn.subtree="ou=users,dc=example,dc=com" by self read

# 允许管理员完全控制
olcAccess: to * by dn="cn=admin,dc=example,dc=com" write

# 允许匿名用户读取公开信息
olcAccess: to dn.subtree="ou=public,dc=example,dc=com" by anonymous read
\end{verbatim}

\chapter{LDAP 安全配置}
\label{chap:ldap_security}

\section{TLS/SSL 加密}
\label{sec:tls_ssl}

为 LDAP 配置 TLS/SSL 加密，保护数据传输安全。

配置步骤：

\begin{verbatim}
# 创建自签名证书
openssl req -newkey rsa:2048 -nodes -keyout /etc/openldap/certs/ldap.key \
  -out /etc/openldap/certs/ldap.csr -subj "/CN=ldap.example.com"
openssl x509 -req -days 365 -in /etc/openldap/certs/ldap.csr \
  -signkey /etc/openldap/certs/ldap.key -out /etc/openldap/certs/ldap.crt

# 设置文件权限
chown ldap:ldap /etc/openldap/certs/ldap.*

# 配置 TLS
cat > /tmp/tls.ldif <<EOF
dn: cn=config
changetype: modify
replace: olcTLSCertificateFile
olcTLSCertificateFile: /etc/openldap/certs/ldap.crt
-
replace: olcTLSCertificateKeyFile
olcTLSCertificateKeyFile: /etc/openldap/certs/ldap.key
EOF

ldapmodify -x -D 'cn=admin,cn=config' -W -f /tmp/tls.ldif
\end{verbatim}

测试 TLS 连接：
\begin{verbatim}
ldapsearch -x -H ldaps://ldap.example.com -b 'dc=example,dc=com' '(objectClass=*)'
\end{verbatim}

\section{密码策略}
\label{sec:password_policy}

配置密码策略，增强安全性。

\begin{verbatim}
# 启用密码策略模块
cat > /tmp/ppolicy.ldif <<EOF
dn: cn=module{0},cn=config
changetype: modify
add: olcModuleLoad
olcModuleLoad: ppolicy.la
EOF

ldapmodify -x -D 'cn=admin,cn=config' -W -f /tmp/ppolicy.ldif

# 配置密码策略
cat > /tmp/ppolicy_config.ldif <<EOF
dn: olcOverlay=ppolicy,olcDatabase={1}mdb,cn=config
objectClass: olcOverlayConfig
objectClass: olcPPolicyConfig
olcOverlay: ppolicy
olcPPolicyDefault: cn=default,ou=ppolicy,dc=example,dc=com
olcPPolicyHashCleartext: TRUE
EOF

ldapmodify -x -D 'cn=admin,cn=config' -W -f /tmp/ppolicy_config.ldif
\end{verbatim}

\section{常见安全威胁}
\label{sec:security_threats}

常见的 LDAP 安全威胁：

\begin{itemize}
    \item \textbf{匿名访问}：未经授权的匿名查询
    \item \textbf{弱密码}：用户使用弱密码
    \item \textbf{中间人攻击}：未加密的通信被窃听
    \item \textbf{注入攻击}：恶意 LDAP 查询注入
    \item \textbf{拒绝服务攻击}：大量请求耗尽服务器资源
\end{itemize}

防护措施：
\begin{itemize}
    \item 配置强访问控制策略
    \item 启用 TLS/SSL 加密
    \item 实施密码策略
    \item 定期备份数据
    \item 监控日志，及时发现异常
\end{itemize}

\chapter{故障排查}
\label{chap:troubleshooting}

\section{常见问题}
\label{sec:common_problems}

配置 LDAP 时常见的问题：

\begin{itemize}
    \item \textbf{连接失败}：检查网络连接、端口、TLS 配置
    \item \textbf{认证失败}：检查用户名、密码、DN 是否正确
    \item \textbf{权限不足}：检查 ACL 配置
    \item \textbf{数据丢失}：检查备份和恢复策略
    \item \textbf{性能问题}：检查索引配置、查询优化
\end{itemize}

\section{日志分析}
\label{sec:log_analysis}

LDAP 日志通常记录在系统日志中。

查看日志：
\begin{verbatim}
# 查看 OpenLDAP 日志
grep slapd /var/log/messages

# 查看 389 Directory Server 日志
grep dirsrv /var/log/messages

# 实时查看日志
tail -f /var/log/messages | grep slapd
\end{verbatim}

\section{诊断工具}
\label{sec:diagnostic_tools}

常用的 LDAP 诊断工具：

\begin{itemize}
    \item \textbf{ldapsearch}：查询目录数据
    \item \textbf{ldapadd}：添加条目
    \item \textbf{ldapmodify}：修改条目
    \item \textbf{ldapdelete}：删除条目
    \item \textbf{ldapwhoami}：验证身份
    \item \textbf{slapcat}：导出目录数据
    \item \textbf{slapadd}：导入目录数据
\end{itemize}

常用命令示例：
\begin{verbatim}
# 验证身份
ldapwhoami -x -D 'cn=张三,ou=users,dc=example,dc=com' -W

# 导出目录数据
slapcat -b 'dc=example,dc=com' > /tmp/backup.ldif

# 导入目录数据
slapadd -b 'dc=example,dc=com' -l /tmp/backup.ldif

# 测试连接
ldapsearch -x -b '' -s base '(objectClass=*)' namingContexts
\end{verbatim}

\chapter{完整操作演练}
\label{chap:practice}

\section{演练环境}
\label{sec:practice_env}

假设环境：
\begin{itemize}
    \item 操作系统：CentOS Stream 9
    \item LDAP 服务器 IP：192.168.1.10
    \item 域名：example.com
    \item 管理员 DN：cn=admin,dc=example,dc=com
\end{itemize}

\section{安装配置 OpenLDAP}
\label{sec:practice_install}

\begin{verbatim}
# 安装 OpenLDAP
sudo dnf install -y openldap openldap-servers openldap-clients

# 启动服务
sudo systemctl start slapd
sudo systemctl enable slapd

# 设置管理员密码
sudo slappasswd

# 创建基础目录结构
cat > /tmp/base.ldif <<EOF
dn: dc=example,dc=com
objectClass: top
objectClass: domain
dc: example

dn: ou=users,dc=example,dc=com
objectClass: top
objectClass: organizationalUnit
ou: users

dn: ou=groups,dc=example,dc=com
objectClass: top
objectClass: organizationalUnit
ou: groups
EOF

# 添加基础目录结构
ldapadd -x -D 'cn=admin,cn=config' -W -f /tmp/base.ldif
\end{verbatim}

\section{配置访问控制}
\label{sec:practice_acl}

\begin{verbatim}
# 创建 ACL 配置
cat > /tmp/acl.ldif <<EOF
dn: olcDatabase={1}mdb,cn=config
changetype: modify
replace: olcAccess
olcAccess: to attrs=userPassword by self write by anonymous auth by * none
olcAccess: to dn.base="" by * read
olcAccess: to * by self write by users read by anonymous auth
EOF

# 应用 ACL 配置
ldapmodify -x -D 'cn=admin,cn=config' -W -f /tmp/acl.ldif
\end{verbatim}

\section{添加用户和组}
\label{sec:practice_users_groups}

\begin{verbatim}
# 添加用户
cat > /tmp/users.ldif <<EOF
dn: cn=张三,ou=users,dc=example,dc=com
objectClass: inetOrgPerson
cn: 张三
sn: 张
givenName: 三
mail: zhangsan@example.com
userPassword: password123

dn: cn=李四,ou=users,dc=example,dc=com
objectClass: inetOrgPerson
cn: 李四
sn: 李
givenName: 四
mail: lisi@example.com
userPassword: password456
EOF

ldapadd -x -D 'cn=admin,dc=example,dc=com' -W -f /tmp/users.ldif

# 添加组
cat > /tmp/groups.ldif <<EOF
dn: cn=开发组,ou=groups,dc=example,dc=com
objectClass: groupOfNames
cn: 开发组
member: cn=张三,ou=users,dc=example,dc=com
member: cn=李四,ou=users,dc=example,dc=com
EOF

ldapadd -x -D 'cn=admin,dc=example,dc=com' -W -f /tmp/groups.ldif
\end{verbatim}

\section{验证配置}
\label{sec:practice_verify}

\begin{verbatim}
# 查询所有用户
ldapsearch -x -b 'ou=users,dc=example,dc=com' '(objectClass=inetOrgPerson)'

# 查询特定用户
ldapsearch -x -b 'ou=users,dc=example,dc=com' '(cn=张三)'

# 查询组成员
ldapsearch -x -b 'ou=groups,dc=example,dc=com' '(cn=开发组)'

# 验证用户身份
ldapwhoami -x -D 'cn=张三,ou=users,dc=example,dc=com' -W
\end{verbatim}

\backmatter

\end{document}